\section{引言}\label{sec:intro}

\subsection{动机}\label{sec:intro:motivation}

自主研究系统有望缩短从问题到实验性答案的距离。核心困难并不只是生成
代码或文本。系统还必须决定尝试什么、判断带噪结果、保存证据、从失败中
恢复,并阻止稿件作出超出运行所证明范围的主张。这些判断构成一种制度:
proposer、evaluator、adversary、writer 与 judge,其行为主要由提示词和
评分规则编码。

固定制度易于复现,却继承初始设计的全部弱点。可演化制度能够适应,但会
产生递归控制问题。因评分变化而获益的组件可能提供改变评分的证据;
reviewer 可能评判自己的后继者;获得更多工具权限的 prompt 可能隐藏暴露
权限扩张的历史。因此,自我改进既是优化问题,也是权限问题。

\subsection{以验证过的主张作为发布条件}\label{sec:intro:verified}

Constitutional ARI-RQGM（自主研究基础设施 ARI 的宪制型 RQGM）从一条
狭义规则出发:模型输出可以\emph{提议}
制度变更,但不能直接\emph{执行}它。确定性内核验证记录形状、hash、
capability、角色分离、转换合法性、audit 连续性、clean-room 输入与
context scope。只有一个 transition engine 拥有 registry 写权限。变更只在
epoch 边界提交；T16 紧急隔离强制关闭当前 epoch，并在同一事务中开启具有
新指纹的 epoch。边界 transaction 要么完整 commit,要么在 replay 时被视为不存在。

固定层还保留最终的结构化数值主张--证据判定。受治理的 writer 与
reviewer 可以比较稿件候选,但不能修改根据已执行证据重算前向声明公式的
verifier。它不是通用自然语言真值检查器。治理可以选择向 gate 提交什么,
却不能放宽 gate。

\subsection{贡献}\label{sec:intro:contrib}

本文作出五项贡献。

\begin{enumerate}
  \item 定义三层权限模型:不可变的 constitutional layer、可演化的
        institutional layer,以及可以提议替代者却不能任命它的 meta layer。
  \item 规定完整的 T1--T21 组件生命周期表,包括 T16 强制紧急边界,以及权限
        矩阵、事件契约和确定性恢复过程。
  \item 闭合从原始产物攻击到 defense、adjudication、validated attack、
        target binding 的问责链；责任可绑定到带纪元来源的研究生成者及
        论文角色，来源含糊时保持无目标。
  \item 通过 selective erasure、clean-room regeneration、replay evaluation
        与 provenance-aware repair 赋予退役以运行含义,其中包括根据保存的
        raw axis 以后继 policy 重新评分。
  \item 把同一宪制应用于 governed paper archive,并报告实现中的 lifecycle
        与 fault injection 证据。
\end{enumerate}

\subsection{本报告的范围}\label{sec:intro:scope}

本文主题是制度协同演化,而非枚举 ARI 的全部执行功能。实验 executor、
domain tool 与 compile pipeline 被视为产生和消费类型化产物的环境,其内部
scheduling strategy 不在讨论范围内。任何遵守记录、来源与 checkpoint
契约的研究 executor 都应满足相同的宪制主张。

我们区分由构造保证的性质与需要实证的主张。registry 的 single-writer
authority、合法转换 topology 以及 fail-closed 状态变更属于结构性质。
本地哈希链能检测中间修改,但不能检测整个保存点和本地固定值一起回退到
旧的合法前缀。新颖性、研究质量、calibration 与 cost effectiveness 的
提升需要对照证据。当前证据是 1,124 项直接相关试验、五个固定故障样例
和一个正常样例;它不能证明一般检测率、归因准确性、研究质量提升或成本
效益。

\subsection{结构}\label{sec:intro:organisation}

\Cref{sec:overview} 定义系统与威胁模型。\Cref{sec:rqgm} 规定
constitutional lifecycle、协同演化、paper governance 与可观测性。
\Cref{chapter:eval} 区分已证明行为和延期主张。\Cref{sec:related} 将本
架构与自修改智能体、自主研究系统及验证工作相比较。\Cref{sec:discussion}
讨论固定宪制的权衡及未解决的扩展问题。附录给出完整核心规格,并逐字收录
RQGM 特有的 runtime prompt seed。
